\documentclass[11pt]{article}
\usepackage{coursenotes}
\begin{document}
\notesheader{11}{Synthetic control}{One treated unit, and a comparison built from many}

\begin{decision}
A change was made in one city, one region, or one flagship store. No single untreated unit looks like it, and picking comparison
units by hand invites picking the ones that flatter the result. Synthetic control builds the comparison by a rule fixed in advance:
a weighted combination of untreated units that tracks the treated unit before the change.
\end{decision}

\begin{goals}
  \item state the synthetic control estimator and explain its constraints;
  \item justify it with a factor model and say why a long, close pre-period fit matters;
  \item design before the post-period: donor pool, fitting window, held-out check;
  \item use in-space and in-time placebos and leave-one-out, and say when a placebo rank is a $p$-value;
  \item recognize its limits (extreme units, spillovers, anticipation) and relate it to difference-in-differences (DiD), synthetic DiD, and matrix completion.
\end{goals}

%=====================================================================================
\sectionone

\subsection{Setup}

Unit $1$ is treated after period $T_0$; units $j = 2, \dots, J+1$ form the \emph{donor pool} and are never treated in the window. Outcomes
$Y_{jt}$ are observed for $t = 1, \dots, T$. Index potential outcomes by adoption date: $Y_{jt}(g)$ is unit $j$'s outcome in period $t$ if it is
first treated in period $g$, and $g = \infty$ means never treated. The target is
$\tau_t = Y_{1t}(T_0 + 1) - Y_{1t}(\infty)$ for $t > T_0$, where $Y_{1t}(T_0 + 1) = Y_{1t}$ is observed and $Y_{1t}(\infty)$ must be constructed.

DiD would compare the treated unit's change with the average donor's. With one treated unit that is often unlike any donor (bigger,
faster-growing, more seasonal), the average donor is a poor counterfactual. If you choose donors by hand instead, you can choose the ones
that flatter the result.

\subsection{The estimator}

\begin{defn}{Synthetic control}{sc}
Choose weights $w = (w_2, \dots, w_{J+1})$ to minimize $\sum_{t=1}^{T_0}\big(Y_{1t} - \sum_j w_jY_{jt}\big)^2$ subject to $w_j \ge 0$ and
$\sum_j w_j = 1$. The synthetic control is $\hat Y_{1t}(\infty) = \sum_j w_jY_{jt}$, and $\hatt_t = Y_{1t} - \hat Y_{1t}(\infty)$ for $t > T_0$.
\end{defn}

Pre-treatment covariates (population, income) can be matched along with pre-period outcomes. Non-negative weights summing to one make the
synthetic unit an \emph{interpolation}: a point inside the donors' range. Unconstrained weights, as in a regression, can fit the pre-period
exactly with combinations that have no real-world analog and behave wildly afterwards. The constraints also make the solution sparse
and explainable, as in ``the synthetic city is 51\% city F, 26\% city E, and 22\% city D''.

\subsection{Why a close pre-period fit is evidence}

\begin{prop}{The factor-model rationale}{factor}
Suppose $Y_{jt}(\infty) = \delta_t + \lambda_t^\top\mu_j + \varepsilon_{jt}$ with common time effects $\delta_t$, unit characteristics $\mu_j$ whose influence
$\lambda_t$ varies over time, and mean-zero noise. If weights match the unit characteristics, $\sum_j w_j\mu_j = \mu_1$, then
$\E[\hat Y_{1t}(\infty)] = \E[Y_{1t}(\infty)]$ for every $t$, including $t > T_0$.
\end{prop}
\begin{proof}
$\sum_j w_jY_{jt}(\infty) = \delta_t + \lambda_t^\top\sum_j w_j\mu_j + \sum_j w_j\varepsilon_{jt}$, using $\sum_j w_j = 1$; with matched characteristics this is
$\delta_t + \lambda_t^\top\mu_1$ plus mean-zero noise.
\end{proof}

The characteristics $\mu_j$ are unobserved, so weights are chosen to match outcomes instead. If a weighted donor combination tracks the treated
unit over a \emph{long} pre-period, during which $\lambda_t$ moves, it must nearly match $\mu_1$, and so keeps tracking it afterwards. DiD is the
special case in which $\lambda_t$ is constant. A close fit over a few periods is easy and uninformative: weights can fit noise.

\subsection{Design before the post-period}

As in Meeting~7, the choices that shape the estimate can be fixed before post-period outcomes are seen.
\begin{enumerate}
  \item \textbf{Donor pool.} Include units similar in the drivers of the outcome; exclude units affected by the treatment (spillovers) or hit by
  their own large shocks.
  \item \textbf{Fitting window and predictors.} Fix them in writing.
  \item \textbf{Held-out check.} Fit on the early pre-period and check that the synthetic unit predicts the late pre-period.
  \item Only then open the post-period.
\end{enumerate}
Weights tuned after seeing the post-period gap are fitted to it.

\subsection{Inference by placebo}

With one treated unit, there is no standard error in the usual sense.

\begin{prop}{When a placebo rank is a $p$-value}{placebo}
Apply the method to every unit in turn as if it were treated (\emph{in-space placebos}) and compute each unit's ratio
$r_j = \text{RMSPE}^{\text{post}}_j/\text{RMSPE}^{\text{pre}}_j$ of root mean squared gaps after and before $T_0$. If the treated unit was chosen at
random among the $J+1$ units, and there is no effect, then the probability that its ratio ranks first is $1/(J+1)$, so rank $k$ gives an exact
$p$-value of $k/(J+1)$.
\end{prop}
\begin{proof}
Under no effect and random choice of the treated unit, the $J+1$ ratios are exchangeable with respect to which unit is labeled treated, so each rank
is equally likely for the treated unit.
\end{proof}

When the treated unit was chosen for a reason (expected to do well, the largest market), the exchangeability fails and the rank is a
\emph{descriptive} measure of how unusual its gap is, not a $p$-value. The ratio discounts placebo units whose pre-period fit was poor; with a
near-perfect pre-period fit the ratio becomes very large or undefined, another reason to read it descriptively. With $J$ donors the smallest
attainable rank-based $p$-value is $1/(J+1)$. \emph{In-time placebos} move the adoption date earlier and check that no gap opens before the true
date. \emph{Leave-one-out} refits without each positive-weight donor; if the estimate depends on one donor, argue for that donor directly.

\subsection{Limits}

\begin{itemize}
  \item \textbf{Extreme treated units.} If the treated unit is the largest or fastest-growing, no convex combination reaches it; the pre-period gap is
  large and the method should not be used as is. Flagship units are exactly this case.
  \item \textbf{Spillovers} to donors (lost orders to the treated city) bias the gap upwards.
  \item \textbf{Anticipation}: date $T_0$ to the announcement if behavior changed then.
  \item \textbf{Local shocks} coinciding with treatment look like effects. Where untreated units share the area but not the treatment, check them
  for a gap at the same date.
\end{itemize}

\subsection{Synthetic control, DiD, and in between}

DiD gives every donor equal weight and allows a constant level difference; synthetic control chooses unequal weights and matches levels. When
the treated unit's path is parallel to the average donor's they agree; when it is not, DiD carries the difference into the estimate and synthetic
control corrects it. Synthetic DiD and matrix completion (Section~3) sit between the two and handle several treated units.

\begin{pitfall}[synthetic-control errors]
Choosing donors after seeing the post-period; a short pre-period; a flagship outside the donors' range; donors that receive spillovers; calling a
placebo rank a $p$-value when the treated unit was chosen for a reason.
\end{pitfall}

%=====================================================================================
\sectiontwo

\paragraph{Setting.} Pujiang Delivery launched same-day grocery delivery in one city (A) at the start of quarter 9. Nine other cities (B to J) did not get it in
the study window. The outcome is monthly orders per 100 residents, averaged over the quarter. The launch is justified if it raises orders by at least
$1.5$ per 100 residents a month within a year. City A was chosen at random among 10 candidate cities by a draw recorded before launch. (The data are
simulated from a factor model, with a planted effect that grows from $1.0$ to $2.2$.)

\paragraph{Step 1: design.} The donor pool is the nine cities, none adjacent to A, and the fitting window is quarters 1--8. For the held-out
check, fitting on quarters 1--6 gives a synthetic city that misses quarters 7--8 by $0.11$ and $0.12$, against a pre-period RMSPE of $0.036$,
which is acceptable.

\paragraph{Step 2: weights and effect.} Fitting on quarters 1--8 gives weights $0.513$ on F, $0.264$ on E, and $0.223$ on D.
\begin{center}
\footnotesize
\setlength{\tabcolsep}{3pt}
\begin{tabular}{lcccccccc|cccc}
\toprule
Quarter & 1 & 2 & 3 & 4 & 5 & 6 & 7 & 8 & 9 & 10 & 11 & 12 \\
\midrule
A (treated) & 14.0 & 16.4 & 16.2 & 16.0 & 18.4 & 20.7 & 20.7 & 20.8 & 23.9 & 26.8 & 27.3 & 27.4 \\
D & 9.0 & 10.4 & 10.4 & 10.6 & 11.9 & 13.1 & 13.1 & 13.6 & 15.0 & 15.7 & 16.3 & 16.4 \\
E & 21.9 & 26.1 & 25.6 & 25.1 & 29.1 & 32.9 & 32.5 & 32.4 & 36.1 & 39.9 & 39.8 & 39.3 \\
F & 12.0 & 13.9 & 13.9 & 13.7 & 15.8 & 17.8 & 17.9 & 17.9 & 19.5 & 21.5 & 21.7 & 21.3 \\
\midrule
Synthetic A & 13.95 & 16.34 & 16.21 & 16.02 & 18.45 & 20.74 & 20.69 & 20.78 & 22.89 & 25.07 & 25.28 & 24.97 \\
Gap & 0.05 & 0.06 & $-0.01$ & $-0.02$ & $-0.05$ & $-0.04$ & 0.01 & 0.02 & 1.01 & 1.73 & 2.02 & 2.43 \\
\bottomrule
\end{tabular}
\end{center}
Check quarter 1: $0.223 \times 9.0 + 0.264 \times 21.9 + 0.513 \times 12.0 = 13.94$ (13.95 with unrounded weights). The gap averages $1.80$ after launch, above $1.5$.

\begin{center}
\begin{tikzpicture}
\begin{axis}[width=0.7\linewidth, height=5cm, xlabel={quarter}, ylabel={orders per 100 residents}, axis lines=left, xmin=1, xmax=12,
  legend style={draw=none, font=\footnotesize, at={(0.02,0.98)}, anchor=north west}]
\addplot[thick, sbgreen] coordinates {(1,14.0) (2,16.4) (3,16.2) (4,16.0) (5,18.4) (6,20.7) (7,20.7) (8,20.8) (9,23.9) (10,26.8) (11,27.3) (12,27.4)};
\addplot[thick, dashed, noteblue] coordinates {(1,13.95) (2,16.34) (3,16.21) (4,16.02) (5,18.45) (6,20.74) (7,20.69) (8,20.78) (9,22.89) (10,25.07) (11,25.28) (12,24.97)};
\draw[extgrey] (axis cs:8.5,12) -- (axis cs:8.5,28);
\legend{city A, synthetic A}
\end{axis}
\end{tikzpicture}

\small Figure 1. City A and its synthetic control; launch after quarter 8.
\end{center}

\paragraph{Step 3: placebos.} Ratios of post- to pre-period RMSPE for every city treated in turn:
\begin{center}
\small
\begin{tabular}{lcccccccccc}
\toprule
City & A & F & G & B & D & E & I & H & J & C \\
\midrule
Pre RMSPE & 0.04 & 0.09 & 0.04 & 0.13 & 1.20 & 2.75 & 0.56 & 0.12 & 0.38 & 0.16 \\
Post RMSPE & 1.87 & 0.74 & 0.09 & 0.22 & 1.67 & 3.66 & 0.58 & 0.11 & 0.29 & 0.12 \\
Ratio & 50.0 & 8.5 & 2.1 & 1.6 & 1.4 & 1.3 & 1.0 & 0.9 & 0.8 & 0.7 \\
\bottomrule
\end{tabular}
\end{center}
A ranks first of 10. Because A was drawn at random among the 10 cities, Result~\ref{prop:placebo} applies and gives $p = 1/10$, the smallest
possible with nine donors.
Had A been picked as the most promising city, the same rank would be descriptive only. Cities E and D, the largest and smallest, cannot be fitted by
the others (pre-RMSPE $2.75$ and $1.20$), so their large post-period gaps mean nothing, and the ratio discounts them. A's ratio is inflated by its very
close pre-fit ($0.04$); read it as ``unusual'', not as a measure of size.

\paragraph{Step 4: robustness.} Dropping F, the largest donor, gives weights on J, D, I, and E, pre-RMSPE $0.067$, and post-period gaps averaging $1.69$.
An in-time placebo at quarter 6 (fitting on quarters 1--5) gives gaps of $-0.19$, $0.06$, and $0.30$ in quarters 6--8, small next to the post-launch gaps.

\paragraph{Step 5: compared with DiD.} An equal-weighted DiD against all nine cities gives $1.46$, below the threshold. A grows faster and swings more
seasonally than the average city, so the average is not a parallel comparison. Here the method changes the decision.

\begin{exercise}[computation]
Verify the synthetic value for quarter 12 from the weights, and the quarter-12 gap.
\end{exercise}
\answer{$0.223 \times 16.4 + 0.264 \times 39.3 + 0.513 \times 21.3 = 3.66 + 10.38 + 10.93 = 24.96$ with the rounded weights ($24.97$ with
unrounded ones); gap $27.4 - 24.97 = 2.43$.}

\begin{exercise}[judgment]
Pujiang Delivery's next launch will be in its largest city, which has twice the order rate of any other. Can synthetic control evaluate it? What would you
propose?
\end{exercise}
\answer{Probably not: no convex combination of smaller cities reaches it, so the pre-period fit will be poor and the gap will mix the effect with the
misfit. Options: match on log outcomes or per-capita rates if the shape, not the level, is what matters; use synthetic DiD (which allows a level shift)
or augmented synthetic control; or randomize the launch within the city (by district, with staggered dates) so the comparison is designed rather
than constructed.}

%=====================================================================================
\sectionthree

\subsection*{Synthetic difference-in-differences}

Synthetic DiD (Arkhangelsky et al.\ 2021) chooses \emph{unit weights} $\omega_j$ (non-negative, summing to one, with a small ridge penalty) so that the
weighted donors' pre-period \emph{trend} matches the treated units' up to a constant, and \emph{time weights} $\lambda_t$ on pre-periods so that the weighted
pre-period resembles the post-period for donors. It then solves the weighted two-way fixed-effects regression
\[
  (\hat\tau, \hat\mu, \hat\alpha, \hat\beta) = \arg\min \sum_{i,t}\big(Y_{it} - \mu - \alpha_i - \beta_t - \tau D_{it}\big)^2\,\hat\omega_i\hat\lambda_t .
\]
Allowing $\alpha_i$ (a level shift) is what synthetic control lacks; using unequal $\omega_i$ is what DiD lacks. It handles several treated units adopting
together, and it remains reliable when either parallel trends or the synthetic-control fit is approximately right. Placebo or jackknife variance
estimators are available.

\subsection*{Matrix completion}

Arrange untreated outcomes in a units-by-periods matrix $\mathbf Y(\infty)$ whose treated cells are missing. Matrix completion (Athey et al.\ 2021) assumes
$\mathbf Y(\infty) = \mathbf L + \boldsymbol\varepsilon$ with $\mathbf L$ of low rank (a factor model with few factors) and estimates
\[
  \hat{\mathbf L} = \arg\min_{\mathbf L} \sum_{(i,t)\ \text{untreated}} (Y_{it} - L_{it})^2 + \lambda\,\|\mathbf L\|_\ast ,
\]
where $\|\cdot\|_\ast$, the sum of singular values, favors low rank. Treated cells' counterfactuals are read from $\hat{\mathbf L}$. Unit and time fixed effects are a
rank-two special case (DiD); synthetic control and horizontal regressions are other special patterns. It uses all the untreated cells, including
treated units' pre-periods and donors' whole histories, and suits staggered adoption.

\begin{extension}[Augmented synthetic control and conformal inference]
When the pre-period fit is imperfect, augmented synthetic control adds an outcome-model correction for the remaining imbalance (Ben-Michael et al.\ 2021). Conformal methods give prediction intervals for the post-period counterfactual under exchangeability of residuals over time.
\end{extension}

\subsection*{Annotated reading}
\begin{readinglist}
\reading{Abadie, Alberto. 2021. ``Using Synthetic Controls: Feasibility, Data Requirements, and Methodological Aspects.''
\emph{Journal of Economic Literature} 59 (2): 391--425. \url{https://doi.org/10.1257/jel.20191450}}{When synthetic control is and is
not appropriate}{All}{1}
\reading{Abadie, Alberto, Alexis Diamond, and Jens Hainmueller. 2010. ``Synthetic Control Methods for Comparative Case Studies:
Estimating the Effect of California's Tobacco Control Program.'' \emph{Journal of the American Statistical Association} 105 (490):
493--505. \url{https://doi.org/10.1198/jasa.2009.ap08746}}{The method and the factor-model argument, with the Proposition 99
example}{Sections 1--3}{2}
\reading{Arkhangelsky, Dmitry, Susan Athey, David A. Hirshberg, Guido W. Imbens, and Stefan Wager. 2021. ``Synthetic
Difference-in-Differences.'' \emph{American Economic Review} 111 (12): 4088--4118.
\url{https://doi.org/10.1257/aer.20190159}}{The synthetic DiD estimator and its properties}{Sections I--III}{3}
\reading{Athey, Susan, Mohsen Bayati, Nikolay Doudchenko, Guido Imbens, and Khashayar Khosravi. 2021. ``Matrix Completion Methods
for Causal Panel Data Models.'' \emph{Journal of the American Statistical Association} 116 (536): 1716--1730.
\url{https://doi.org/10.1080/01621459.2021.1891924}}{Low-rank counterfactuals as a unifying view}{Sections 1--4}{3}
\reading{Ben-Michael, Eli, Avi Feller, and Jesse Rothstein. 2021. ``The Augmented Synthetic Control Method.'' \emph{Journal of the
American Statistical Association} 116 (536): 1789--1803. \url{https://doi.org/10.1080/01621459.2021.1929245}}{Correcting imperfect
fit}{Sections 1--3}{3}
\end{readinglist}

\end{document}
